% ══════════════════════════════════════════════════════════════
%  Section 1: Introduction
% ══════════════════════════════════════════════════════════════

\section{Introduction}\label{sec:introduction}

Two traditions compete for the task of predicting asset returns. Statistical machine-learning methods---neural networks, random forests, gradient-boosted trees---estimate flexible functions of observable characteristics and deliver strong out-of-sample (OOS) performance, yet they treat economic theory as irrelevant input \citep{gu2020empirical}. Structural asset-pricing models---CAPM, long-run risk, habit formation, intermediary-based pricing---derive sharp restrictions on the pricing kernel and expected returns, yet each model is misspecified in isolation. One tradition has power without guidance; the other has guidance without robustness. Neither alone is satisfactory.

We propose a framework that draws on both. Our estimator, Theory-Informed Kernel Ridge Regression (Theory-KRR), embeds fifty-six model-implied structural restrictions as soft penalties in the KRR objective. Each restriction $j$ enters with its own penalty multiplier~$\mu_j$, and the entire vector $\bm{\mu} = (\mu_1, \dots, \mu_J)$ is selected by cross-validation. The mechanism is automatic: restrictions from misspecified theories receive~$\mu_j \to 0$---the data shut them down---while restrictions from empirically useful theories receive~$\mu_j > 0$, steering the estimator toward economically sensible regions of the reproducing kernel Hilbert space~$\Hcal$. No single structural model needs to be correct. The data adjudicate.

The restrictions we impose are not generic shape constraints such as monotonicity or convexity \citep{bryzgalova2023asset}. They are genuine model-implied restrictions: Euler-equation moment conditions, factor-loading constraints, demand-system equilibrium conditions, and stochastic discount factor bounds, each derived from a specific economic theory. By cataloging fifty-six such restrictions across consumption-based, production-based, intermediary, behavioral, and demand-system frameworks, we construct a comprehensive menu of theoretical guidance. Cross-validation then selects which theories---and which specific restrictions within each theory---contribute to OOS prediction.

A methodological parallel clarifies the logic. The LLM-Lasso of \citet{bybee2025llmlasso} replaces the agnostic $\ell_1$ penalty in Lasso with heterogeneous, text-informed penalties: large-language-model predictions set a prior over which covariates matter, and cross-validation adjusts a single scalar that scales the text-derived penalty vector. Our approach applies the same principle in a different domain. Where LLM-Lasso uses text priors to inform penalization in \emph{parameter space}, Theory-KRR uses economic-theory priors to inform penalization in \emph{function space}. The mathematical structure is analogous: both replace a single, uninformed regularization parameter with a vector of informed, heterogeneous penalties, and both let cross-validation arbitrate between the prior and the data. Theory-KRR is, in this sense, the function-space generalization of LLM-Lasso.

What does this framework deliver? Theory-KRR outperforms standard KRR and several benchmark ML methods in OOS return prediction, measured by $R^2_{\text{OOS}}$ and portfolio Sharpe ratios. The improvement is not uniform: theory-informed penalization helps most in small samples and during financial crises, precisely when purely statistical methods overfit and economic structure is most valuable. Beyond predictive gains, the cross-validated multipliers~$\bm{\mu}$ produce an empirical ranking of economic theories by their marginal OOS forecasting value. Consumption-based restrictions \citep{campbellcochrane1999,bansalyaron2004} and intermediary-based restrictions \citep{hekrishnamurthy2013,adrianetulamuir2014} receive the largest multipliers, indicating that these theories provide the most useful guidance for return prediction. Production-based and behavioral restrictions contribute modestly. Several models receive~$\mu_j$ indistinguishable from zero, consistent with those theories being empirically irrelevant for forecasting in our sample.

This paper makes three contributions. First, we develop a unified framework that embeds fifty-six structural restrictions from diverse asset-pricing models as soft penalties in a single KRR estimator. The framework nests standard KRR (all~$\mu_j = 0$) and hard-constrained estimation (any~$\mu_j \to \infty$) as special cases, and it accommodates restrictions of heterogeneous functional form---moment conditions, inequality constraints, and functional equations---within a single objective. Second, we establish a formal connection between theory-informed function-space penalization and the LLM-Lasso approach to parameter-space penalization. This connection unifies two strands of the literature---structural econometrics and text-based ML---under a common principle of informed regularization. Third, we produce the first empirical ranking of economic theories by their marginal contribution to OOS forecasting, providing a new lens through which to evaluate the empirical content of competing asset-pricing models.

Our work speaks to several strands of the literature. The ML asset-pricing literature, initiated by \citet{gu2020empirical}, demonstrates that flexible nonparametric methods substantially outperform linear models in predicting the cross-section of returns. Subsequent work explores neural-network architectures \citep{chen2024deep}, autoencoder factor models with no-arbitrage constraints \citep{Gu2021autoencoder}, conditional factor models \citep{kelly2019characteristics}, and Gaussian process regression applied to equity returns \citep{Filipovic2022gaussian}. \citet{Kelly2024virtue} provide a theoretical foundation for this success, showing that complex, overparameterized models are not penalized by benign overfitting and systematically outperform simpler ones. These methods are powerful but agnostic about economic structure. Our contribution is to show that injecting economic theory into the regularization of a kernel method yields both predictive gains and economic interpretability.

A growing body of work recognizes that theory and statistics are complements, not substitutes, for return prediction. The paper closest to ours in spirit is \citet{Chen2026teaching}, who show that pre-training neural networks on synthetic data generated by structural models---and then fine-tuning on empirical data---substantially improves option-pricing accuracy, even when the structural model is misspecified. Our approach shares the same fundamental insight---misspecified theory still helps---but differs in three important respects. First, we use soft penalization in the objective function rather than pre-training: structural content enters directly as penalty terms with cross-validated weights, not as an initialization from synthetic data. Second, we target the cross-section of equity returns rather than option pricing. Third, the convexity of our KRR objective guarantees a global optimum and yields tractable asymptotics, whereas neural-network pre-training is inherently non-convex. Most importantly, our penalty weights provide a continuous measure of each theory's marginal contribution to prediction, whereas transfer learning offers a binary use-or-not decision for each structural model.

A second closely related paper is \citet{Bryzgalova2023bayesian}, who develop a Bayesian framework that averages over 2.25~quadrillion linear asset-pricing models formed by subsets of 51~observable factors. Their identification-strength-based priors impose heterogeneous shrinkage---factors that are weakly identified receive stronger shrinkage---and the resulting Bayesian model-averaging stochastic discount factor dominates individual factor models out of sample. Like us, they produce an empirical ranking of economic models. The key differences are the objects being ranked and the estimation framework. Bryzgalova, Huang, and Julliard rank observable factor portfolios within a linear SDF; we rank structural restrictions from named theoretical models within a nonparametric RKHS. Their priors are calibrated from factor identification strength; our penalty weights are selected by temporal cross-validation, tying the ranking directly to out-of-sample predictive performance.

More broadly, the question of which factors or models survive in a high-dimensional setting has been explored from several angles. \citet{Kozak2020shrinking} propose economically motivated shrinkage of SDF coefficients proportional to the eigenvalues of characteristic-based principal components, motivated by no-arbitrage bounds on Sharpe ratios. \citet{Feng2020taming} develop a double-selection LASSO procedure to test whether newly proposed factors have marginal explanatory power beyond existing ones. \citet{Harvey2016cross} document severe data snooping in factor discovery and argue for higher significance thresholds. \citet{Barillas2018comparing} and \citet{Kan2013pricing} provide formal Bayesian and frequentist frameworks for pairwise model comparison. Our paper complements this literature by moving from factor-level selection in linear models to theory-level selection in a nonparametric estimator, where the ``theories'' are entire families of structural restrictions rather than individual tradable factors.

The econometrics literature on constrained nonparametric estimation provides our mathematical foundation. \citet{babii2024structural} develop the theory of penalized kernel regression with functional constraints, establishing convergence rates and conditions under which constraints improve estimation. \citet{gabaix2024asset} show how demand-system equilibrium conditions can be imposed on asset-pricing estimators, and \citet{bryzgalova2023asset} impose shape restrictions (monotonicity, convexity) on pricing-kernel estimators. In the penalized GMM literature, \citet{Caner2009lasso} and \citet{Caner2014adaptive} develop LASSO-type and elastic-net GMM estimators with oracle properties for moment selection. \citet{Freyberger2020dissecting} use adaptive group LASSO to select characteristics nonparametrically. Our work extends this line by moving from reduced-form shape restrictions to model-implied structural restrictions, and by allowing the data to select among competing restrictions via cross-validation rather than imposing all restrictions with equal weight. The representer theorem, generalized to our multi-penalty setting following \citet{Scholkopf2002learning} and the foundational results of \citet{Kimeldorf1971correspondence} and \citet{Wahba1990spline}, ensures that the solution lies in the span of kernel evaluations even with the added structural penalties.

The paper proceeds as follows. Section~\ref{sec:framework} presents the mathematical framework, defining the Theory-KRR objective and establishing its properties. Section~\ref{sec:restrictions} catalogs the fifty-six structural restrictions we draw from the asset-pricing literature, organized by model class. Section~\ref{sec:estimation} addresses practical estimation, including cross-validation of the multiplier vector~$\bm{\mu}$, computational strategies, and implementation details. Section~\ref{sec:empirical} describes the data and baseline models. Section~\ref{sec:results} presents the main empirical results: OOS predictive performance, the theory ranking, and robustness exercises. Section~\ref{sec:conclusion} concludes.
